\section{Remark 5}\label{s:remark}

Remark~5 has four sentences. The first says what the theorem answers, the second and third what it does not, and the fourth where the missing parts are to be found. The remark is the only place in EXI-01-a where the text speaks about itself, and it does so with precision: each of its claims about what the theorem does and does not answer can be checked against the definition, the assumptions and the proof, and each is correct.

\lem{9}{The theorem answers the question of what the meaning of life is, and of whether meaning is available to everyone.}

\lem{9.1}{The theorem answers the question of what the meaning of life is}

\paragraph{The theorem or the definition?} Strictly, the question what the meaning of life is is answered by Definition~1, not by Theorem~4. The theorem says that a life of positive meaning is available to every agent; it does not say what meaning is. The sentence is best read with \enquote{the theorem} understood as the result together with the definition in whose terms it is stated, as one speaks of a theorem of arithmetic as answering a question that its definitions make precise. The alternative reading, on which the remark claims that the theorem alone answers the first question, would make the remark false, and the principle of charity forbids it where another reading is available.

\paragraph{What kind of answer.} The answer to \enquote{what is the meaning of life?} that Definition~1 gives is a real definition of the property in virtue of which lives are meaningful. The comment on \Lm{1} has argued that this is the second of the three questions the title can raise. The remark confirms the reading: it treats the question of what the meaning of life is as answered, and the only answer in the text is Definition~1. A reader who asked the first question, what human life as such is for, will find the remark's claim surprising, and the text owes her the explanation, given in the comment on \Lm{1}, that the second question does not depend on the first. EXI-11 and EXI-06 support the explanation from two sides: meaning does not require a cosmic purpose (Theorem~3.3 of EXI-11-a), and the universe's indifference takes nothing from it (EXI-06-b, its Proposition~2.1).

\lem{9.2}{and of whether meaning is available to everyone}

The second question is answered by Theorem~4, with \enquote{everyone} read as every agent capable of intentional action (\Lm{4.2}). The answer is yes. The phrase \enquote{available to everyone} is the phrase of the family conclusion, and the remark thus identifies Theorem~4 as the text's statement of the second half of that conclusion. The comment on \Lm{5.3} has argued that the availability is best read as indexed to times, so that meaning remains available to an agent for as long as she is one.

\lem{10}{It does not answer the question of whether any particular life is meaningful, which depends on what is done with it.}

\paragraph{Availability is not possession.} The sentence draws the distinction on which EXI-01-d builds its account of what the availability claim does and does not imply: from availability one cannot infer actuality. Whether a particular life is meaningful depends on whether it in fact consists in active engagement with projects of objective worth, and Theorem~4 says only that some life available to its subject does.

\paragraph{\enquote{what is done with it}.} The relative clause makes the meaning of a particular life depend on what is done with it. The passive voice leaves open who does it, and the omission is apt. Much of what is done with a life is done by its subject: whether she engages with what is worthwhile is largely up to her. But some is done by others and by circumstance: the range of projects within her reach, the time and strength she has for them, the education that lets her see their worth. A voluntarist reading, on which the meaning of a life is wholly up to its subject, would overstate the case; a fatalist reading, on which it is wholly up to circumstance, would contradict Theorem~4. The passive voice accommodates both contributions, and EXI-01-d states how they are divided: the possibility of meaning is secured whatever circumstances do, and its degree depends on what the subject does and on the conditions under which she does it (its Proposition~2.4).

\paragraph{The left-to-right direction.} The sentence uses Definition~1 in the direction the proof does not use: a life in which nothing worthwhile is engaged with is not meaningful. The comment on \Lm{2.3} noted that the definition is a biconditional and that the remark depends on this. The Lean module, which formalizes only the positive case of Definition~1, cannot represent this sentence (Section~\ref{s:lean}).

\paragraph{The first-person question.} The question that Remark~5 sets aside, whether a particular life is meaningful, is the question a person most often asks about her own life, and Definition~1 says how it is to be answered. It divides into two questions that can be asked separately. The first is whether, and how far, her life consists in active engagement: whether there are things she cares about in a way that issues in sustained, directed activity. This is a question about her, and she is in a good position to answer it, though not an infallible one, since a person can mistake habit for care and care for habit. The second is whether the objects of her engagement have worth that does not depend on her caring about them. This is a question about the objects, and she is in no better position to answer it than anyone else who knows the facts and reflects well. The two questions have different epistemologies, and the division explains a familiar feature of crises of meaning: they take two forms. In one, the person finds that she no longer cares, as Mill did; the remedy, if there is one, lies in her. In the other, she finds that what she cared about did not deserve it, as Ivan Ilyich did; the remedy lies in finding something that does. Definition~1 gives each form its own diagnosis.

\paragraph{The comparative question.} Remark~5 says that whether a particular life is meaningful depends on what is done with it. The same holds for how meaningful it is, and for whether one life is more meaningful than another. The comment on \Lm{2.3} has argued that the degree depends on at least three variables, the share or amount of the life given to engagement, the depth of the engagement and the worth of its objects, and that some comparisons may be indeterminate. A particular life can therefore be made more meaningful in three ways: by giving more of it to engagement, by engaging more deeply with what one already engages with, or by turning to objects of greater worth. EXI-11-a's Principle~P2 states the third; EXI-17-a's remark that achievement may add to the degree of meaning (its Remark~3.5) bears on the second, since success can deepen engagement and enlarge the worth realized through it. The first is the one most within a person's control and the one most limited by circumstance, and it is the one on which EXI-01-d's Proposition~2.4 bears most directly.

\lem{11}{Nor does it answer the sceptical challenge that, \emph{sub specie aeternitatis}, nothing matters \citep{Nagel1971}: Assumption 2 denies the premise of that challenge rather than refuting it.}

This is the most important sentence of the remark and the one whose interpretation is most disputed in the discipline. It makes two claims. The first is that the theorem does not answer the challenge. The second, after the colon, says why: Assumption~2 denies the premise of the challenge but does not refute it. We take the phrases in turn and then give the analysis that the sentence requires.

\lem{11.1}{the sceptical challenge}

\paragraph{\enquote{sceptical}.} The adjective places the challenge in the family of sceptical arguments, and it does so on Nagel's own authority. \citet{Nagel1971} compared the sense of the absurd to epistemological scepticism: as the sceptic about knowledge steps back from our beliefs and finds that every justification we can offer depends on assumptions we cannot justify without circularity, so the reflective subject steps back from her concerns and finds that every justification of them depends on further concerns open to the same question. In both cases, Nagel argued, we cannot refute the sceptic and we cannot stop believing, or caring; and in both cases the right response is neither to pretend to a refutation nor to abandon what cannot be justified, but to live with the discrepancy.

\paragraph{\enquote{challenge}.} The noun is chosen with care. A challenge is not a thesis. The sceptic about knowledge need not assert that we know nothing; she asks how we know, and finds every answer wanting. The sceptic about meaning need not assert that nothing matters; she asks what justifies our taking anything to matter, and finds every answer open to the same question. The distinction matters for what Assumption~2 can do against the challenge. Against a thesis that nothing matters, Assumption~2 is a contradictory thesis. Against a challenge, it is an answer that the challenger will not accept.

\lem{11.2}{\emph{sub specie aeternitatis}}

The phrase is Spinoza's: the mind conceives things under a species of eternity when it conceives them as following from the necessity of the divine nature, and not in relation to a time and place (\emph{Ethics} V, propositions 29--31). Nagel adopts it for the standpoint from which the reflective subject regards her life as one among others, without presupposing the importance of anything in it, and he observes that it is a standpoint we can occupy without leaving our lives, since it is the standpoint of reflection itself \citep{Nagel1971}. In \emph{The View from Nowhere} he gave the standpoint a general theory, as the objective standpoint that we reach by abstracting from the particularities of our own point of view, and returned in its last chapter to the meaning of life \citep{Nagel1986}. EXI-04-a defines the corresponding notion as the detached standpoint: the standpoint from which an agent considers her ends as the ends of one creature among others, without presupposing their importance, and asks what justifies her concern with them (its Definition~2.2).

The Spinozan phrase carries an implication that Nagel's use drops. For Spinoza, to conceive things under the species of eternity is to conceive them truly, and the highest virtue of the mind lies in knowledge of this kind (\emph{Ethics} V, proposition 25). For Nagel, the standpoint discloses not the truth about our concerns but the absence of a vindication of them. EXI-01-a uses the phrase in Nagel's sense, and its citation of Nagel confirms this.

\lem{11.3}{nothing matters}

\paragraph{The formula.} \enquote{Nothing matters} is the slogan of the challenge in its strongest form. It has a history in the analytic literature before Nagel. \citet{Hare1972} described a young visitor who, after reading Camus's \emph{L'\'Etranger}, came to the conclusion that nothing matters and was distressed by it, and argued that the conclusion rests on a confusion about the word \enquote{matters}: to say that something matters is to express concern about it, and since the visitor was concerned about many things, the conclusion that nothing matters to him was false, and the conclusion that nothing matters \emph{simpliciter} was not something he could make sense of. Hare's answer is expressivist, and it is not the answer EXI-01-a gives, since on Assumption~2 whether something has worth is not settled by whether anyone is concerned about it. But Hare's diagnosis, that the challenge trades on an ambiguity between \enquote{matters to someone} and \enquote{matters}, survives the change of metaethics, and the analysis below uses it.

\paragraph{\enquote{matters} and \enquote{worth}.} The challenge says that nothing \emph{matters}; Assumption~2 says that some projects have \emph{worth}. The words are not synonyms, but in the context of the challenge they are equivalent: to matter, \emph{sub specie aeternitatis}, is to have importance that does not depend on its mattering to anyone in particular, and that is objective worth in the sense of Assumption~2 or a stronger one. The restriction of Proposition~\ref{p:restricted} does not affect the equivalence, since the detached standpoint abstracts from the attitudes of the subject who occupies it, and worth that does not depend on the attitudes of those engaged is worth that the detached subject has not presupposed by abstracting from her own engagement.

\lem{11.4}{Assumption 2 denies the premise of that challenge rather than refuting it}

\paragraph{Which premise?} The sentence speaks of \enquote{the premise} of the challenge, in the singular, and does not say which premise it is. A challenge, as the comment on \Lm{11.1} noted, need not have premises in the way an argument does. But a challenge that is to have the force of an objection to Assumption~2 must be reconstructible as an argument for the negation of something Assumption~2 asserts. The question is which of its premises Assumption~2 denies.

We reconstruct the challenge as follows.

\begin{enumerate}[label=(C\arabic*)]
\item If anything matters, its mattering can be vindicated from the detached standpoint: there is a justification of it that the detached standpoint is compelled to accept.
\item Nothing's mattering can be vindicated from the detached standpoint: every justification of concern is open to the question what justifies it, and no answer is forced on that standpoint.
\item Therefore nothing matters.
\end{enumerate}

(C2) is Nagel's central claim and the content of Lemma~3.2 of EXI-04-a, which proves it for any reflective agent: if the agent answers that a consideration has objective worth, the detached standpoint can register the claim but cannot be compelled to endorse it, since what she would exhibit is her own conviction, which the step back has placed in question. (C1) is the requirement that mattering be vindicable from the detached standpoint. The argument is valid.

\begin{proposition}[The premise denied]\label{p:vind}
Given (C2), Assumption~2 entails the negation of (C1). The premise of the sceptical challenge that Assumption~2 denies is (C1), the requirement that whatever matters be vindicable from the detached standpoint.
\end{proposition}

\begin{proof}
Assumption~2 entails that something has objective worth, and so, by the equivalence noted at \Lm{11.3}, that something matters; this contradicts (C3). Since the argument is valid, Assumption~2 is inconsistent with the conjunction of (C1) and (C2). The text accepts (C2): it is not mentioned in EXI-01-a, but it is proved in EXI-04-a, whose conclusion the family accepts, and Remark~5 does not dispute it. So the premise Assumption~2 denies is (C1).
\end{proof}

\paragraph{Denying and refuting.} A premise is refuted when its falsity is derived from premises that the one who holds it accepts, or that she is rationally required to accept. It is denied when its negation is asserted, or entailed by what is asserted, without such a derivation. Assumption~2 denies (C1): it entails, together with (C2), that (C1) is false, but it does not derive the falsity of (C1) from anything the sceptic accepts, since the sceptic does not accept Assumption~2. Remark~5 is therefore exact. The theorem does not answer the challenge, because the challenge calls into question a premise of the theorem's interpretation; and Assumption~2 denies a premise of the challenge without refuting it, because the denial rests on Assumption~2 itself.

\paragraph{Is the denial question-begging?} A denial that rests on the very claim the challenge calls into question may seem to beg the question. It does not, for two reasons. First, the denial is not offered as an answer to the sceptic; Remark~5 says expressly that the theorem does not answer her. It is offered as a statement of where the text stands, so that the reader can see that the text's conclusion is conditional on rejecting (C1), and can assess (C1) for herself. Second, (C1) has a cost that can be assessed independently of Assumption~2. It is a vindication requirement: it says that a claim of worth is true only if it can be vindicated from a standpoint that presupposes no worth. The analogous requirement in epistemology, that a belief is knowledge only if it can be vindicated from a standpoint that presupposes no belief, leads to scepticism about everything, including the beliefs on which the sceptical argument itself relies, and Nagel's analogy is an analogy in exactly this respect. A requirement that, applied in the one domain, would eliminate all knowledge, and that has no better support in the other, may be rejected without a refutation. MTE-05 makes a parallel point about the argument from queerness: an argument against objective value that would equally show mathematics false proves too much.

\paragraph{The testimony of EXI-04-a.} EXI-04-a gives the same reading of Remark~5. In its reply to the objection that objective worth dissolves the absurd, it says that seriousness is justified if the projects are worthwhile, that the absurd survives because that justification is not one the detached standpoint must accept, and that its definition of the absurd concerns the availability of a forced justification, not the existence of a sound one; and it concludes that this is the sense in which, as the remark at the end of EXI-01-a puts it, the assumption of objective worth denies the premise of the sceptical challenge without refuting it. In the terms of Proposition~\ref{p:vind}: the existence of a sound justification is what Assumption~2 asserts; the availability of a forced one is what (C1) requires and (C2) denies. EXI-04-a and EXI-01-a agree that (C2) is true and (C1) is false, and they agree that the falsity of (C1) is asserted rather than proved.

\subsection{The absurd}\label{s:absurd}

Remark~5 says that a discussion of the absurd is given in the Commentary (\Lm{12}). The analysis of \Lm{11.4} has shown where EXI-01-a stands on Nagel's challenge. We now relate that position to the principal treatments of the absurd and state what follows for meaning.

\paragraph{Nagel.} \citet{Nagel1971} rejected the standard arguments for the absurdity of life: that our lives are short, that we are small in the universe, that what we do will not matter in a million years, that every justification ends somewhere. Each, he argued, either fails to show what it claims or presupposes the sense of absurdity it purports to explain: if our lives are absurd at seventy years, they would be no less absurd if they lasted through eternity; if nothing we do now will matter in a million years, then nothing that will happen in a million years matters now. The real source of the absurd is the collision between the seriousness with which we live and the detachment of which we are capable. His conclusion was that if, \emph{sub specie aeternitatis}, there is no reason to believe that anything matters, then that does not matter either, and we can approach our lives with irony rather than with heroism or despair. EXI-04 takes over the conclusion and gives it a proof: the absurd is ineliminable for any reflective agent (Theorem~3.3 of EXI-04-a), despair and revolt do not fit it (Propositions~3.4 and~3.5), and irony does (Theorem~3.6).

EXI-01-a is consistent with this conclusion, and Proposition~\ref{p:vind} shows exactly how. Nagel's argument establishes (C2), which EXI-01-a accepts. It does not establish (C1), and Nagel's own conclusion does not need it: Nagel says that the detached standpoint gives no reason to believe that anything matters, not that it gives a reason to believe that nothing does. The sceptic who infers (C3) has added (C1), and it is (C1) that Assumption~2 denies. The absurd, in Nagel's sense, and meaning, in the sense of Definition~1, are therefore compatible.

\begin{proposition}[Meaning and absurdity]\label{p:compat}
Every life of positive meaning lived by a reflective agent is absurd in the sense of EXI-04-a, and its absurdity does not reduce its meaning.
\end{proposition}

\begin{proof}
By Theorem~3.3 of EXI-04-a, the life of every reflective agent capable of intentional action is absurd, whatever its content; so a meaningful life of such an agent is absurd. Absurdity consists in the absence of a forced justification of seriousness from the detached standpoint (Definition~2.3 of EXI-04-a). The degree of meaning of a life, by Definition~1, depends on the extent to which it consists in active engagement with projects of objective worth, and neither the extent of engagement nor the worth of its objects depends on whether that worth can be vindicated from the detached standpoint, since Assumption~2, which fixes what worth is, makes no reference to vindication. So absurdity does not reduce meaning.
\end{proof}

\paragraph{Camus.} \citet{Camus1955} located the absurd in the confrontation between the human longing for clarity and meaning and the unreasonable silence of the world, and he took three responses into view: physical suicide, philosophical suicide, the leap of faith by which one escapes the absurd into a meaning beyond it, and revolt, the lucid refusal either to escape or to submit. His Sisyphus, conscious of the futility of his labor and scornful of the gods who imposed it, is the figure of revolt, and Camus concludes that one must imagine him happy. EXI-04-a argues that revolt does not fit the absurd, because it presents the world as an adversary, and an indifferent world is not one (its Proposition~3.5 and Remark~3.7; on indifference, EXI-06-a). EXI-04-b distinguishes the absurdity of a life from the futility of a task: Sisyphus's task is futile, both because it cannot bring about its outcome and because the outcome has no value (its Definition~1.1), while the absurdity of a life is compatible with its projects' succeeding and having worth (its Proposition~1.2).

The distinction is essential to EXI-01-a. Camus's Sisyphus is a figure of the absurd and of futility together, and the image draws its power from the futility. On Definition~1, his life is not meaningful, however lucid and however happy, because the object of his labor has no worth; EXI-04-b makes the same point and notes that Taylor's Sisyphus is in the same position. The human situation, on the analysis of EXI-04, is absurd without being futile: our projects can succeed, and many of their objects have worth. The meaning that Definition~1 describes is meaning that Camus's Sisyphus lacks and that the absurd does not take from us. What remains of Camus's conclusion, EXI-04-b argues, is the claim that full engagement with one's ends is compatible with lucidity about the absence of their vindication (its Proposition~2.1), and that is the ironic response that Proposition~\ref{p:compat} shows to be compatible with meaning.

\paragraph{Tolstoy.} Tolstoy's crisis \citep{Tolstoy1987} has a different structure from Nagel's and Camus's. It is driven not by detachment but by death: the thought that everything he did would come to nothing, since he would die and his works would be forgotten, made his goods taste of nothing, as the honey in the fable no longer tasted sweet to the traveller who saw the dragon and the mice. His resolution was religious: he found in the faith of ordinary people a meaning that death does not destroy, because it relates the finite life to the infinite.

Tolstoy's argument is that what comes to nothing was nothing. EXI-11-a states it as an objection and replies that the inference from transience to worthlessness is invalid: a friendship that ends with death was a friendship, and the relief of a patient's pain was relief while it lasted (its reply to Objection~1). EXI-19-b adds that the limit death sets is not a source of meaning and does not make anything engaged with more or less worthwhile (its Proposition~1.1). On Definition~1, the meaning of a life consists in what the life consists in, and the comment on \Lm{2.4} has drawn the consequence that it is fixed by the life's course and is not undone by what happens after it. Tolstoy's premise, that worth realized at a time is cancelled by later loss, is therefore denied by the constitutive reading of Definition~1, and the denial is supported by the theory of value: on ETH-01, the value of an outcome is fixed by idealized endorsement of it, and idealized endorsement of the relief of pain does not wait on whether the relief is remembered.

There remains something in Tolstoy's case that this reply does not touch. His goods lost their taste: the engagement that had filled his life failed. On Definition~1, his life during the crisis lost meaning, not because its objects lost their worth, but because he could no longer engage with them. This is the structure of Mill's crisis (Section~\ref{s:defence}, D3), and it shows that the thought of death can threaten meaning through the subjective component, by undermining engagement, even if it cannot threaten it through the objective component. EXI-05 holds that fear of death is not rational, though regret at its prospect may be. If the fear is not fitting, its power to undermine engagement is the power of an unfitting attitude, and the remedy is the clarity about what death does and does not take that EXI-02, EXI-05 and EXI-19 together supply. Volume~III, Part~C, develops the point.

\paragraph{Hume's remedy.} Nagel's analogy between the absurd and scepticism has a precedent that EXI-04-a considers as an objection. Hume, at the end of the first book of the \emph{Treatise} (1.4.7), describes how the sceptical reflections that had reduced him to melancholy were dispelled when he dined, played backgammon and conversed with his friends, and how, returning to them afterward, he found them cold and strained. The objection is that forgetting, not irony, is the appropriate response to the absurd. EXI-04-a replies that forgetting is not a response to the absurd but an interruption of the reflection that discloses it, and that the question is what attitude fits the absurd when one attends to it (its reply to Objection~4). For EXI-01-a, Hume's case has a further interest. Hume's return to ordinary life is a return to engagement, with friends, with games and with conversation, and his melancholy was a failure of engagement induced by reflection. The case therefore illustrates, as Tolstoy's and Mill's do, the way in which reflection can threaten meaning through its subjective component without touching its objective one. What reflection showed Hume was that his beliefs could not be vindicated from the standpoint of the sceptic; it did not show that they were false, and his friends, his backgammon and his conversation were no less worth engaging with for his having, for a time, been unable to engage with them.

\paragraph{The point of view of the universe.} Sidgwick held that the good of any one individual is of no more importance, from the point of view of the universe, than the good of any other \citep[III.xiii]{Sidgwick1907}. \citet{Williams1982} criticized the phrase, arguing that there is no such point of view and that ethics cannot take its standards from it. The criticism bears on (C1). If there is no point of view of the universe from which our concerns could be vindicated, the demand that they be so vindicated is not a demand that could be met by anything, and its failure to be met shows nothing about whether anything matters. The detached standpoint of EXI-04-a is not the point of view of the universe; it is a standpoint the reflective subject can occupy. But (C1) requires vindication from that standpoint in its most detached form, and Williams's criticism suggests that the most detached form is not a standpoint from which anything could be vindicated or condemned.

\paragraph{Cosmic insignificance.} \citet{Kahane2014} distinguished the claim that we are insignificant from the cosmic standpoint from the claim that our lives lack meaning, and argued that the first, if true, depends on the claim that we are surrounded by much greater value elsewhere in the universe, not on our smallness. EXI-06 holds that the universe does not care about us, and EXI-06-b argues that its indifference takes nothing from meaning. Both are consistent with Definition~1, which makes meaning depend on the worth of what we engage with, not on how that worth compares with the worth of things elsewhere. A small contribution to a great good is engagement with something of objective worth, and its meaning is not reduced by the existence of greater goods, though its significance, in Kahane's sense, may be.

\paragraph{Summary.} The absurd, in Nagel's sense, is real and ineliminable, and EXI-01-a accepts it. It does not show that nothing matters, because that conclusion requires (C1), which Assumption~2 denies. It does not reduce the meaning of a life, because meaning depends on engagement and worth, not on vindication (Proposition~\ref{p:compat}). It is distinct from futility, which does remove meaning, and Camus's image confuses the two. It is distinct from the threat of death, which can undermine engagement but not worth, and Tolstoy's argument from transience to worthlessness is invalid. The fitting response to it is irony, as EXI-04 holds, and irony is compatible with the full engagement on which meaning depends. Volume~II, Part~C, gives the history of the absurd at length and develops the analysis of (C1).

\lem{12}{A defence of Assumption 2, and a discussion of the absurd, are given in the commentary (EXI-01-b).}

The last sentence refers to the present work. It makes two commitments, and both are met in this volume: the defence of Assumption~2 is Section~\ref{s:defence}, and the discussion of the absurd is the analysis of \Lm{11.4} together with Section~\ref{s:absurd}. Both are extended in Volume~II.

The sentence also tells the reader something about the text. It is a text that knows it is incomplete and says where its completion is to be found. The pattern is that of axiomatic texts with commentaries, in which the axioms are stated without defence and the defence is given separately. It shows that the brevity of EXI-01-a is a choice of form, not a judgment that its assumptions need no defence: the text separates the derivation, which fits on one page, from the defence of its premises, which does not.
